\documentclass[10pt]{beamer}

\usepackage{appendixnumberbeamer}
\usepackage{booktabs}
\usepackage[scale=2]{ccicons}
\usepackage{pgfplots}
\usepackage{xspace}
\usepackage{bookmark}
\usepackage{amssymb}
\usepackage{mathtools}
\usepackage[normalem]{ulem}
\usepackage[T1]{fontenc}
\usepackage[sfdefault,book]{FiraSans}
\usepackage{FiraMono}
\usepackage{fontawesome}
\usepackage{hyperref}
\usepackage{listings}

\renewcommand{\thefootnote}{\fnsymbol{footnote}}
\renewcommand{\MintedPygmentize}{/home/thiago/.local/bin/pygmentize}

\definecolor{light-gray}{gray}{0.95}
\renewcommand\lstlistingname{Código}
\DeclareCaptionFormat{listing} {
  \parbox{\textwidth}{\hspace{-0.2cm}#1#2#3}
}
\DeclareCaptionFont{black}{\color{black}}
\captionsetup[lstlisting]{
  format=listing,
  labelfont=black,
  textfont=black,
  singlelinecheck=true,
  margin=0pt,
  font={tt,footnotesize,bf}
}

\lstset{
  basicstyle=\footnotesize\ttfamily,
  escapeinside={\%*}{*)},
  mathescape=true,
  showspaces=false,
  showtabs=false,
  showstringspaces=false,%
  rulesepcolor=\color{black},
  frame=shadowbox,
  upquote=true,
  literate=
  {á}{{\'a}}1 {é}{{\'e}}1 {í}{{\'i}}1 {ó}{{\'o}}1 {ú}{{\'u}}1
  {Á}{{\'A}}1 {É}{{\'E}}1 {Í}{{\'I}}1 {Ó}{{\'O}}1 {Ú}{{\'U}}1
  {à}{{\`a}}1 {è}{{\`e}}1 {ì}{{\`i}}1 {ò}{{\`o}}1 {ù}{{\`u}}1
  {À}{{\`A}}1 {È}{{\'E}}1 {Ì}{{\`I}}1 {Ò}{{\`O}}1 {Ù}{{\`U}}1
  {ä}{{\"a}}1 {ë}{{\"e}}1 {ï}{{\"i}}1 {ö}{{\"o}}1 {ü}{{\"u}}1
  {Ä}{{\"A}}1 {Ë}{{\"E}}1 {Ï}{{\"I}}1 {Ö}{{\"O}}1 {Ü}{{\"U}}1
  {â}{{\^a}}1 {ê}{{\^e}}1 {î}{{\^i}}1 {ô}{{\^o}}1 {û}{{\^u}}1
  {Â}{{\^A}}1 {Ê}{{\^E}}1 {Î}{{\^I}}1 {Ô}{{\^O}}1 {Û}{{\^U}}1
  {Ã}{{\~A}}1 {ã}{{\~a}}1 {Õ}{{\~O}}1 {õ}{{\~o}}1
  {œ}{{\oe}}1 {Œ}{{\OE}}1 {æ}{{\ae}}1 {Æ}{{\AE}}1 {ß}{{\ss}}1
  {ű}{{\H{u}}}1 {Ű}{{\H{U}}}1 {ő}{{\H{o}}}1 {Ő}{{\H{O}}}1
  {ç}{{\c c}}1 {Ç}{{\c C}}1 {ø}{{\o}}1 {å}{{\r a}}1 {Å}{{\r A}}1
  {€}{{\euro}}1 {£}{{\pounds}}1 {«}{{\guillemotleft}}1
  {»}{{\guillemotright}}1 {ñ}{{\~n}}1 {Ñ}{{\~N}}1 {¿}{{?`}}1
}

\setminted{
  fontsize=\footnotesize,
  style=bw,
  frame=single,
  labelposition=topline,
}

\pretitle{\begin{center}\normalsize\bfseries\MakeUppercase{Programação 2 -- }\MakeUppercase}
\author{\normalsize Prof. Thiago Cavalcante}
\date{\vspace{-4ex}}

\geometry{
  top=0.0cm,
  left=1.0cm,
  right=1.0cm,
  bottom=1.0cm
}

\setlength\columnsep{30pt}

\usetikzlibrary{arrows}

\makeatletter
\AddEnumerateCounter{\PaddingUp}{\two@digits}{A00}
\AddEnumerateCounter{\PaddingDown}{\two@digits}{A00}
\newcommand\PaddingUp[1]{\expandafter\two@digits\csname c@#1\endcsname}
\newcommand\PaddingDown[1]{\PaddingUp{#1}\addtocounter{#1}{-2}}
\makeatother

\makeatletter
\def\enumalphalphcnt#1{\expandafter\@enumalphalphcnt\csname c@#1\endcsname}
\def\@enumalphalphcnt#1{\AlphAlph{#1}}
\makeatother
\AddEnumerateCounter{\enumalphalphcnt}{\@enumalphalphcnt}{aa}

\pagenumbering{gobble}

\title{Programação 2}
\author{Thiago Cavalcante  -- thiago.cavalcante@penedo.ufal.br}
\institute{Universidade Federal de Alagoas -- UFAL \\ Campus Arapiraca \\ Unidade de Ensino de Penedo}
\titlegraphic{\hfill\includegraphics[height=1.5cm]{images/brasao-ufal.eps}}


\subtitle{Aula 4}
\date{13 de novembro de 2019}

\begin{document}

\maketitle

\begin{frame}{Ponteiros}
  \huge
  \textbf{Ponteiros}
  \vfill
  \LARGE
  Tipo \textbf{especial} de variável que armazena um \textbf{endereço na memória} ao invés de um valor
  \vfill
  \large
  \setlength{\leftmargini}{0pt}
  \begin{itemize}
      \item [] Toda informação está armazenada na memória
      \item [] Ponteiros "apontam" para endereços de memória
      \item [] programador $\rightarrow$ nome, programa $\rightarrow$ endereço
  \end{itemize}
\end{frame}

\begin{frame}{Ponteiros}
  \huge
  \textbf{Ponteiros podem ser \alert{perigosos \faExclamationTriangle}}
  \vfill
  \large
  Podem apontar para um espaço na memória que está sendo usado para um outro propósito
\end{frame}

\begin{frame}[fragile]{Ponteiros}
  \huge
  \textbf{Declaração de ponteiro}
  \vfill
  \Large
  \begin{verbatim}
tipo_do_ponteiro *nome_do_ponteiro;
  \end{verbatim}
  \vfill
  \large
  Não confundir com o operador de multiplicação!
\end{frame}

\begin{frame}{Ponteiros}
  \huge
  \textbf{Inicialização e atribuição}
  \vfill
  \large
  \begin{itemize}
    \item \texttt{int *p;} (aponta para um \textbf{lugar indefinido} \faExclamationTriangle)
    \item \texttt{p = NULL;} (aponta para \textbf{lugar nenhum} -- diferente)
    \item \texttt{p = }\&\texttt{x;} (aponta para variável \texttt{x})
    \item \texttt{printf("\%d", *p)} (imprime o valor guardado na variável apontada por \texttt{p})
    \item \texttt{int *p2 = p;} (ponteiro \texttt{p2} aponta para o mesmo lugar que \texttt{p} \faExclamationTriangle)
  \end{itemize}
  \vfill
  \large
  \&: operador de endereçamento
\end{frame}

\begin{frame}{Ponteiros}
  \huge
  \textbf{Operadores}
  \vfill
  \Large
  \textbf{*}: Usado na \textbf{declaração} de um ponteiro e na \textbf{obtenção do conteúdo} para onde ele aponta (\alert{dois usos distintos})
  \vfill
  \Large
  \textbf{\&}: Usado para \textbf{obter o endereço} de uma variável
\end{frame}

\begin{frame}{Ponteiros}
  \huge
  \textbf{Aritmética com ponteiros}
  \vfill
  \Large
  \begin{itemize}
    \item Adição e subtração: incrementam ou decrementam o tamanho \textbf{em bytes} do tipo do ponteiro \faExclamationTriangle
    \item Outras operações: não são realizadas no ponteiro, mas no seu conteúdo (\textbf{\texttt{*p}})
  \end{itemize}

\end{frame}

\begin{frame}{Ponteiros}
  \huge
  Também é possível comparar ponteiros com operadores relacionais: \texttt{==}, \texttt{!=}, \texttt{>}, \texttt{<}, \texttt{>=}, \texttt{<=}

\end{frame}

\begin{frame}{Ponteiros}
  \huge
  \textbf{Ponteiros genéricos}
  \vfill
  \large
  \begin{itemize}
    \item \texttt{void *nome\_do\_ponteiro};
    \item Aponta para qualquer tipo de dado \faExclamationTriangle
    \item Acesso precisa ser feito com um \textit{typecast}: \texttt{*(int*)p}
    \item Incremento é feito de byte em byte
  \end{itemize}
\end{frame}

\begin{frame}[fragile]{Ponteiros}
  \huge
  \textbf{Ponteiros e arrays}
  \vfill
  \Large
  O nome de um array é apenas um \textbf{ponteiro para o primeiro elemento} do array
  \vfill
  \large
  \begin{lstlisting}
ponteiro = nome_array
ponteiro = %*\textnormal{\&}*)nome_array[0]

ponteiro[i] = 2;
*(ponteiro + i) = 2;
  \end{lstlisting}
\end{frame}

\begin{frame}[fragile]{Ponteiros}
  \huge
  \textbf{Array de ponteiros}
  \vfill
  \Large
  \begin{lstlisting}
tipo_de_dado *nome_do_array[tamanho];

nome_do_array[indice] = %*\textnormal{\&}*)variavel
  \end{lstlisting}
\end{frame}

\begin{frame}[fragile]{Ponteiros}
  \huge
  \textbf{Ponteiro para ponteiro}
  \vfill
  \Large
  \begin{verbatim}
tipo_de_dado **nome_do_ponteiro;
  \end{verbatim}
\end{frame}
\end{document}
